\documentclass[12pt]{article}

\usepackage[a4paper,margin=24mm]{geometry}
\usepackage{amsmath,amssymb,bm}
\usepackage{booktabs}
\usepackage{graphicx}
\usepackage{microtype}
\usepackage[numbers,sort&compress]{natbib}
\usepackage[hidelinks]{hyperref}
\usepackage{xcolor}

% The generated result macros retain the historical V15 names. TeX normally
% terminates a control word before a digit, so temporarily classify 1 and 5 as
% letters while reading the generated file and access the names through
% \csname below. This keeps the frozen reporting output reproducible.
\def\restoreVfifteendigitcatcodes{\catcode`1=12\catcode`5=12}
\catcode`5=11
\catcode`1=11
% Generated from V15 aggregate results; do not edit manually.
\newcommand{\V15Trajectories}{48}
\newcommand{\V15Decisions}{34,560}
\newcommand{\V15ClustersPerModel}{6}
\newcommand{\V15GPUHours}{49.41}
\newcommand{\V15OriginalGPUHours}{19.19}
\newcommand{\V15ReconstructionGPUHours}{49.41}
\newcommand{\V15OrphanFormalDecisions}{197}
\newcommand{\V15OrphanFormalGPUHours}{0.328}
\newcommand{\V15UnrecordedInfrastructureCalls}{1}
\newcommand{\V15HOne}{42.263}
\newcommand{\V15HOneCI}{24.316 to 59.303}
\newcommand{\V15HTwo}{0.05438}
\newcommand{\V15HTwoCI}{0.03459 to 0.07548}
\newcommand{\V15HThree}{0.04845}
\newcommand{\V15HThreeCI}{0.02190 to 0.07526}
\newcommand{\V15HFour}{52.541}
\newcommand{\V15HFourCI}{35.406 to 68.673}
\newcommand{\V15HOneDisposition}{supported}
\newcommand{\V15HTwoDisposition}{supported}
\newcommand{\V15HThreeDisposition}{supported}
\newcommand{\V15HFourDisposition}{supported}
\newcommand{\V15QwenCorrQuench}{0.00238}
\newcommand{\V15QwenCorrQuenchCI}{-0.01750 to 0.02277}
\newcommand{\V15GraniteCorrQuench}{-0.02032}
\newcommand{\V15GraniteCorrQuenchCI}{-0.05486 to 0.00947}
\newcommand{\V15QwenTauMemory}{-2.46995}
\newcommand{\V15QwenTauMemoryCI}{-5.79214 to 0.17724}
\newcommand{\V15GraniteTauMemory}{-4.14109}
\newcommand{\V15GraniteTauMemoryCI}{-19.88286 to 11.60067}
\newcommand{\V15QwenBinderQuench}{-4.80504}
\newcommand{\V15QwenBinderQuenchCI}{-7.42824 to -2.20832}
\newcommand{\V15GraniteBinderQuench}{0.16643}
\newcommand{\V15GraniteBinderQuenchCI}{-0.09447 to 0.52063}

\restoreVfifteendigitcatcodes
\newcommand{\E}{\mathbb E}
\newcommand{\KL}{D_{\mathrm{KL}}}
\newcommand{\figroot}{../../results/collective_agent_statmech_v15/figures/pdf}
\newcommand{\resultfigure}[4]{%
\begin{figure}[t]\centering
\IfFileExists{\figroot/#1}{\includegraphics[width=#2]{\figroot/#1}}{\fbox{\parbox[c][45mm][c]{0.88\textwidth}{\centering Figure generated after formal analysis:\\[2mm]\texttt{\detokenize{#1}}}}}
\caption{#3}\label{#4}\end{figure}}

\title{Statistical-mechanical characterization of memory and quench response in state-separated LLM-agent networks}
\author{Anonymous working manuscript}
\date{23 August 2026}

\begin{document}
\maketitle

\begin{abstract}
We treat a network of state-separated, locally informed large-language-model
(LLM) agent instances as a finite interacting stochastic system.  Each
persistent identity owns a categorical belief, a committed action, bounded
private memory, a local field, an inbox and outbox, and typed local authority;
a random-sequential scheduler offers update opportunities but does not choose
the resulting state.  We distinguish the complete augmented simulator state,
its observable belief--action projection, and a rolling macroscopic
representation.  This separation exposes how omitted private memory can act
as a hidden historical coordinate.  We report collective order, uncertainty,
dependence, an effective symmetric-layer energy, fluctuations, and
coarse-grained path-reversal divergence.  We first correct a derived-analysis
error in an earlier frozen Qwen quench study without changing any trajectory.
We then conduct a prospectively frozen two-family experiment with
\csname V15Trajectories\endcsname{} trajectories,
\csname V15Decisions\endcsname{} local decisions, and six
independent graph/environment clusters per model.  A field reversal and
restoration is paired with nominal evolution, genuine persistent memory, and a
format-matched scrambled-history control.  The independent-family quench
effect is \csname V15HOne\endcsname{} (95\% interval
\csname V15HOneCI\endcsname) within-model distance units.  Genuine
memory changes bias-adjusted reversal divergence relative to Markovized and
scrambled-history controls by \csname V15HTwo\endcsname{}
(\csname V15HTwoCI\endcsname) and \csname V15HThree\endcsname{}
(\csname V15HThreeCI\endcsname) nats per attempted update.  The fixed
early-to-late restoration contrast is \csname V15HFour\endcsname{}
(\csname V15HFourCI\endcsname) distance units.  All four effects meet their
frozen criteria, although the memory contrasts are larger and more uniform in
Granite than in Qwen.  An out-of-sample kinetic
surrogate identifies which response components admit a simple local closure
and which remain model- and history-dependent.  These finite-size results
support a reduced statistical-mechanical description without interpreting
decoding temperature as physical temperature, reference energy as literal
energy, or projected path divergence as exact thermodynamic entropy
production.
\end{abstract}

\noindent\textbf{Keywords:} statistical mechanics of agents; language-model
agents; quench dynamics; entropy rate; multi-information; temporal asymmetry.

\section{Introduction}

Networks of language-model agent instances are interacting many-component
systems with unusually complicated local transition rules.  A component can
retain private history, condition on a natural-language observation, receive
only selected messages, and execute a typed local action.  Even when model
weights are fixed, the collective trajectory depends on graph structure,
asynchronous scheduling, message content, bounded memory, and decoding
randomness.  Population averages alone can hide this organization.  The
statistical-mechanics question is whether carefully defined reduced variables
can connect local response to collective order, fluctuation, persistence, and
driven reorganization.

The kinetic Ising model provides the canonical example in which stochastic
local updates generate collective relaxation \cite{glauber1963}.  Related
logit dynamics has an equilibrium stationary law when interactions derive
from a potential \cite{blume1993}.  Nonreciprocal interactions can instead
produce collective structures and responses unavailable to reciprocal
potential systems \cite{fruchart2021,zhang2023}.  State-space currents and trajectory
reversal quantify time asymmetry in Markov processes
\cite{schnakenberg1976,seifert2005,seifert2012,roldan2012}.  Those exact
constructions are references rather than assumptions here.  An observable
LLM belief--action process may fail first-order closure because a prompt can
contain persistent private history and semantically structured packets.
Coarse-graining can hide entropy production or generate estimator-dependent
lower bounds \cite{kawaguchi2013,busiello2019,kutvonen2015}.  We therefore use
the more precise term \emph{coarse-grained pathwise irreversibility} for our
empirical block-reversal statistic.

Recent studies have begun to apply statistical physics to interacting LLM
populations.  Group size can change basins and collective misalignment
\cite{flint2025}; topology and self/social weighting alter conformity
\cite{han2026}; intrinsic bias can dominate fitted neighbor coupling
\cite{denobili2026}; debate noise produces finite-size ordering crossovers
\cite{okawa2026}; and large populations can be summarized by a small family
of fitted dynamical regimes \cite{el2026}.  Our contribution is complementary.
We retain separate belief and committed-action layers, explicit local
authority, inbox and outbox boundaries, random-sequential updates, persistent
private memory, a Markovized comparison, and a controlled quench followed by
a counter-quench.  We place direct Qwen \cite{qwen2024} and Granite
\cite{granite33model} trajectories beside an out-of-sample kinetic closure.

The central claim is descriptive and explanatory, not performance-oriented.
Statistical-mechanical variables provide a compact language for uncertainty,
dependence, compatibility, fluctuation, and temporal response that is not
contained in a single magnetization.  No task reward, controller improvement,
human response, or operational benefit is evaluated.  Nor do two pinned
models establish universality.

\resultfigure{figure01_augmented_state_architecture.pdf}{0.96\textwidth}{%
Agent information boundaries and the augmented-to-macroscopic map.  The
complete state $\Xi_t$ includes private memory, graph and schedule state; the
evaluator observes $Y_t=\phi(\Xi_t)$ and constructs rolling
$Z_t=\psi(Y_{t-w+1:t})$.  The scheduler offers a turn and delivers a packet but
does not supply the model's belief, action, memory, or signal.}{fig:architecture}

\section{Stochastic process and information boundaries}

\subsection{Local agent state}

For agent $i$ at attempted update $t$, the primary categorical coordinates are
belief $b_i(t)\in\{-1,+1\}$ and committed action
$a_i(t)\in\{-1,+1\}$.  Its bounded local state also includes confidence,
commitment status, workload, a memory code, at most three private history
entries, an inbox, and an outbox.  Display labels are counterbalanced and
decoded back to latent spins.  The microscopic projection is
\begin{equation}
Y_t=(\bm b_t,\bm a_t,\bm c_t,\bm m_t),
\end{equation}
but this is not the complete simulator state.

Let
\begin{equation}
\Xi_t=(S_{1,t},\ldots,S_{N,t},G_t,{\cal D}_t,q_t,R_t)
\end{equation}
denote all agent-private states $S_{i,t}$, graph state $G_t$, delivery operator
${\cal D}_t$, quench phase $q_t$, and the explicitly specified randomness
state $R_t$.  With fixed model weights and tokenizer, reconstructed local
prompts, and a defined scheduler, $\Xi_t$ induces a time-inhomogeneous
transition kernel under the external quench protocol.  The evaluator's map
$Y_t=\phi(\Xi_t)$ deliberately omits some prompt history and randomness.
Consequently $Y_t$ and any rolling macrostate need not be Markov.

At each attempted update, a random permutation gives every agent exactly one
opportunity per sweep.  Only the scheduled instance receives its private
observation, current local state, currently delivered inbox, and, in a memory
condition, its own bounded history.  It returns a validated typed object
containing belief, action, confidence, commitment, memory state, outgoing
signal, and tool action.  One bounded repair may correct syntax; an invalid
result causes no centrally selected scientific action.  A serialized packet
is sent on one graph edge.  Fingerprints of nonrecipient peer-private states
are compared around every transition.

\subsection{Markovized, persistent, and scrambled history}

The Markovized prompt is reconstructed from the explicitly recorded current
state and does not carry free conversational history.  The persistent prompt
adds the instance's last three own entries in the fixed format
\texttt{time, belief, action, memory}.  Thus two trajectories with identical
visible $(\bm b_t,\bm a_t)$ can have distinct conditional future laws:
\begin{equation}
P(Y_{t+1}\mid Y_t,M_t)\ne P(Y_{t+1}\mid Y_t,M'_t).
\end{equation}
If $M_t$ is projected out, conditional dependence on older visible states can
remain.  This is a concrete mechanism for increased block reversal divergence,
not a proof of physical dissipation.

The V15 scrambled-history control isolates prompt presence and approximate
length.  It uses only the same agent's past opportunity times, while belief,
action, and memory content are deterministically randomized from a frozen
control seed.  No donor agent, future time, or hidden environmental outcome is
used.  Genuine and scrambled entries have the same section, count, field
names, and format.  Token counts and their paired differences are measured;
turn-by-turn equality is not assumed.

\section{Collective observables}

\subsection{Order and effective compatibility}

Belief and action order and their overlap are
\begin{equation}
m_b=\frac{1}{N}\sum_i b_i,\qquad
m_a=\frac{1}{N}\sum_i a_i,\qquad
q_{ba}=\frac{1}{N}\sum_i b_i a_i.
\end{equation}
Disagreement is the fraction of connected pairs with unequal states.  For a
symmetric comparison layer $W_s=(W+W^\top)/2$, we define
\begin{equation}
H_{\mathrm{ref}}=-\frac{J_b}{2}\bm b^\top W_s\bm b
-\frac{J_a}{2}\bm a^\top D_s\bm a
-K\bm a^\top\bm b-\bm h^\top\bm b-\bm g^\top\bm a .
\label{eq:href}
\end{equation}
We report $e_{\mathrm{ref}}=H_{\mathrm{ref}}/N$.  Equation
(\ref{eq:href}) is an effective compatibility coordinate anchored in a
symmetric reference model.  We do not infer a Gibbs law for the LLM process
and do not call this literal energy.  Likewise
$N\operatorname{Var}(e_{\mathrm{ref}})$ is an energy-fluctuation observable,
not heat capacity unless an independently justified effective temperature
exists.

The belief susceptibility is reported without physical-temperature
normalization,
\begin{equation}
\chi_b=N\operatorname{Var}(m_b),
\end{equation}
within a prescribed rolling window.  Three secondary observables resolve
spatial organization, temporal persistence, and the shape of the finite-size
order-parameter distribution.  First, for a phase window $W$ we define the
connected belief correlation at graph distance $d$,
\begin{equation}
C_b(d;W)=\frac{1}{|{\cal P}_d|}\sum_{(i,j)\in{\cal P}_d}
\left[\langle b_i b_j\rangle_W-\langle b_i\rangle_W
\,\langle b_j\rangle_W\right],\qquad
{\cal P}_d=\{(i,j):i<j,\ d_G(i,j)=d\}.
\label{eq:connected-correlation}
\end{equation}
Shortest paths use the actual reciprocal modular delivery support.  The
bracket averages post-update observations at attempted-update resolution
within one phase of one
trajectory; node pairs are averaged inside that trajectory and never counted
as independent replicates.

Second, the normalized magnetization autocorrelation and its fixed truncation
are
\begin{align}
\rho_m(\ell)&=\frac{\langle[m_b(t)-\bar m_b]
\,[m_b(t+\ell)-\bar m_b]\rangle}
{\langle[m_b(t)-\bar m_b]^2\rangle},\nonumber\\
\tau_{\mathrm{int}}^{(\ell_{\max})}&=\frac{1}{2}+
\sum_{\ell=1}^{\ell_{\max}}\rho_m(\ell).
\label{eq:autocorrelation}
\end{align}
The primary descriptive truncation is $\ell_{\max}=2N=32$ attempted updates;
$N$ and $3N$ are fixed sensitivities.  Zero-variance phase series are left
undefined rather than assigned artificial persistence.  The quantity is a
finite-window correlation sum reported in attempted-update units, not evidence
of critical slowing down or an effective sample-size estimate.  Its absolute
short-lag scale includes persistence induced by changing at most one agent on
each random-sequential update; matched arms share that update convention.
Quench and restoration windows can also be nonstationary, so their sums are
finite-window persistence diagnostics rather than equilibrium integrated
correlation times.

Third, we calculate the Binder cumulant
\begin{equation}
U_4=1-\frac{\langle m_b^4\rangle}
{3\langle m_b^2\rangle^2}
\label{eq:binder}
\end{equation}
separately for each phase and trajectory \cite{binder1981}.  Near-zero second moments are
reported as undefined.  With only $N=16$, equation (\ref{eq:binder}) is a
finite-size distribution-shape diagnostic; no crossing, critical point, or
thermodynamic-limit transition is inferred.  A post-reconstruction estimator
sensitivity repeats the calculation over each full phase and its early and
late halves.  It compares the mean of the six trajectory-level cumulants with
a pooled-moment construction; both uncertainty calculations resample complete
graph--environment clusters, never updates.

\subsection{Entropy and dependence}

For discretized collective words $x$ in a rolling window, Shannon
configuration
entropy is
\begin{equation}
S_{\mathrm{config}}=-\sum_x \hat p(x)\log\hat p(x).
\end{equation}
This is the usual plug-in Shannon functional \cite{shannon1948}.
The entropy rate is a bounded-history plug-in estimate of next-word
unpredictability.  Single-variable marginal entropy measures local occupancy;
pairwise mutual information measures dependence between variables.  Total
correlation \cite{watanabe1960},
\begin{equation}
{\cal T}=\sum_{k=1}^{2N}H(X_k)-H(X_1,\ldots,X_{2N}),
\end{equation}
separates marginal uncertainty from synchronous joint dependence.

The plug-in joint entropy of 32 binary variables in a short window has a large
finite-sample floor \cite{paninski2003}.  We therefore retain the raw estimate
but independently circular-shift each variable 200 times.  Shifts preserve
each marginal occupancy and within-variable temporal ordering while breaking
synchronous cross-variable relations.  We report raw, null mean, their
untruncated difference, and a normalized difference divided by the sum of
marginal entropies.  The same audit is applied to pairwise and edge mutual
information.  A negative adjusted estimate remains negative.

\subsection{Rolling macrostate}

The primary five-sweep representation is
\begin{align}
Z_t=(&m_b,m_a,q_{ba},e_{\mathrm{ref}},v_E,S_{\mathrm{config}},h_{\mathrm{rate}},
{\cal T}_{\mathrm{adj}},I_{\mathrm{pair,adj}},I_{\mathrm{edge,adj}},\nonumber\\
&\chi_b,\rho_b,D_b)_t .
\end{align}
Three- and seven-sweep windows are frozen sensitivities.  This is an
observable reduced representation, not a complete thermodynamic state.

\section{Pathwise temporal asymmetry}

For a discrete observed word sequence and block length $L$, let
$\hat P_L(x_0,\ldots,x_L)$ be the regularized empirical forward block law.
We compute
\begin{equation}
{\cal I}_L=\frac{1}{L}
\KL\left[\hat P_L(x_0,\ldots,x_L)\parallel
\hat P_L(x_L,\ldots,x_0)\right]
\label{eq:pathkl}
\end{equation}
in nats per attempted update.  The primary $L=3$ estimator uses pseudocount
0.5.  Values at $L\in\{2,4\}$ and pseudocounts 0.1 and 1.0 are sensitivities.
Five hundred time permutations provide a panel-specific finite-length floor,
and the adjusted statistic is observed (\ref{eq:pathkl}) minus its mean floor.
Matched arms share the same shuffle tape.  The machine-readable sensitivity
table retains both the empirical permutation-null interval and the Monte Carlo
standard error of its mean.  These quantify different objects: dispersion of
the null statistic and numerical uncertainty in the estimated bias floor,
respectively.  The added uncertainty columns reproduce, rather than replace,
the frozen floor mean and adjusted contrast.

For a fully observed stationary Markov chain with transition kernel $P$ and
stationary law $\pi$, the current
\begin{equation}
J_{xy}=\pi_xP_{xy}-\pi_yP_{yx}
\end{equation}
can define an exact Markov entropy-production rate under the usual support
conditions.  We do not apply that interpretation to projected LLM trajectories.
Persistent histories, semantic content, and bounded observations can violate
first-order closure.  Equation (\ref{eq:pathkl}) is instead a coarse-grained
time-reversal diagnostic and, under additional assumptions, may be a lower
bound on hidden irreversibility \cite{roldan2012,busiello2019,liang2023}.

\resultfigure{figure02_memory_discovery_replication.pdf}{0.94\textwidth}{%
Discovery, prospective replication, and V15 cross-model memory contrasts are
shown separately.  Points are graph/environment-cluster effects; intervals
resample complete clusters.  V12 and V13 are not retroactively pooled into the
V15 confirmatory interval.}{fig:memory}

\section{Quench protocol and nominal geometry}

Each trajectory has 15 baseline sweeps, 15 intervention sweeps, and 15
restoration sweeps.  In the field-quench arm the vector of balanced private
fields changes sign at the first boundary and returns to its initial value at
the second.  Nominal trajectories retain the initial field.  This follows the
logic of a finite-system quench protocol, while making no quantum or
thermodynamic-limit analogy.  The initial
change is the quench; restoration is the counter-quench.

For each held-out graph/environment cluster and model, the nominal manifold is
fitted only from nominal trajectories in the other five clusters.  Training
medians impute nonfinite coordinates; training means and scales standardize
features; shrinkage covariance with frozen ridge fraction defines a
Mahalanobis-like distance.  The nominal 95th-percentile threshold is computed
from those training distances.  The held-out cluster contributes neither a
scale nor a threshold.

We summarize peak departure, integrated response, path length, final residual,
and threshold re-entry.  The prospective V15 recovery contrast is
\begin{equation}
\Delta d_{\mathrm{rec}}=
\overline d_{31:35}-\overline d_{41:45},
\end{equation}
which can be positive, zero, or negative.  It replaces a historical maximum
minus final-window quantity whose sign was structurally constrained.

\resultfigure{figure03_v14_quench_time_series.pdf}{0.94\textwidth}{%
Corrected V14 time series for nominal and field-reversal trajectories.  The
yellow interval marks the quench; the dotted boundary marks restoration.
Reference energy is effective, entropy is configuration diversity, and
distance is measured in the cluster-excluded nominal geometry.}{fig:v14series}

\resultfigure{figure04_v14_cluster_recovery.pdf}{0.94\textwidth}{%
All V14 field-reversal cluster trajectories, including the high-response
cluster.  The derived audit replaces the held-out-panel threshold with the
training-cluster nominal threshold.  The historical H3 sign statistic is not
used as directional evidence.}{fig:v14recovery}

\section{Prospective V15 design}

The original model is Qwen/Qwen2.5-7B-Instruct at immutable revision
{\footnotesize\texttt{a09a3545\allowbreak 8c702b33\allowbreak eeacc393\allowbreak d1030632\allowbreak 34e8bc28}}.  The independent family is
IBM Granite-3.3-8B-Instruct at revision
{\footnotesize\texttt{51dd4bc2\allowbreak ade4059a\allowbreak 6bd87649\allowbreak d68aa11e\allowbreak 4fb2529b}}.  Both run with Transformers
4.55.4, PyTorch 2.8.0+cu128, bitsandbytes 0.47.0, NF4 double quantization, BF16
computation, decoding temperature 0.5, top-$p=0.9$, and at most 96 generated
tokens.  The model's decoding temperature is not an inferred thermodynamic
temperature.

Each family contributes six new graph/environment clusters.  Each cluster has
$N=16$, a reciprocal modular fixed-degree graph, coupling $J=0.8$, balanced
disordered initialization, and four matched trajectories: nominal Markovized,
field-quench Markovized, field-quench persistent memory, and field-quench
scrambled history.  Graphs, fields, initial states, label maps, update orders,
recipient variates, and inference seeds are paired where technically
meaningful.  A complete graph/environment trajectory cluster is the
inferential unit.

The engineering pilot is restricted to loading, schema validity, occupancy,
both transition directions, timing, privacy, delivery, prompt-token balance,
latency, and projected compute.  It does not calculate a quench, memory, or
irreversibility contrast.  The formal design contains
\csname V15Trajectories\endcsname{} trajectories and
\csname V15Decisions\endcsname{} attempted local decisions.  Complete
unfavorable trajectories are retained.

The hypotheses are frozen as follows.  H1 compares Granite field-quench peak
departure with matched nominal evolution at one-sided $\alpha=0.02$.  H2
compares persistent with Markovized adjusted reversal divergence.  H3 compares
persistent with scrambled history.  H4 tests the fixed recovery contrast.
H2--H4 form a one-sided Holm family with total $\alpha=0.03$.  Exact cluster
sign flips are the primary tests; deterministic 10,000 cluster bootstraps give
intervals.  The sign-flip enumeration is exact conditional on the observed
paired magnitudes under a sign-symmetry null; it is not described as treatment-
label randomization.  Agent and time-point counts do not increase the
inferential sample size.

\section{V14 correction and audit}

The V14 audit preserves every frozen prompt, completion, categorical choice,
and trajectory.  The original source computed a field panel's recovery
threshold from its own baseline, despite protocol text specifying a
training-only threshold.  Audit version 1.1 passes an explicit cluster-to-
training-threshold map from leave-one-cluster-out fitting into recovery
summaries.  A versioned archive retains the committed historical JSON, tables,
README, claims matrix, and macros.

The frozen H3 number was the maximum recovery distance minus the final-five
mean.  It remains numerically reproducible, but this construction is
nonnegative for nearly every nonconstant trajectory.  Its raw and Holm-adjusted
sign-flip probabilities are retained as historical audit values and removed
from inferential support.  The corrected machine-readable fields distinguish
whether the numerical criterion was met, whether a directional test is valid,
and whether the time-resolved path is consistent with return.

The audit also executes the omitted 10,000 cluster-preserving label
permutations.  All four panels stay together within each cluster, and every
replicate refits imputation, scaling, logistic regression, and leave-one-
cluster-out predictions.  Three-, five-, and seven-sweep nominal geometries
are fitted independently.  Each coordinate and observable family is deleted
in turn.  Dependence estimates receive circular-shift null floors.  These are
delayed prespecified sensitivities, not a prospective V14 experiment.

\section{Results: field-quench replication}

The Granite-family H1 contrast is \csname V15HOne\endcsname{} distance units
(95\% interval \csname V15HOneCI\endcsname); its frozen disposition is
\csname V15HOneDisposition\endcsname{}.  Figure
\ref{fig:crossquench} displays model-specific distances rather than pooling
their absolute scales.  A positive contrast indicates a field-driven
departure beyond nominal variation; a nonpositive or uncertain contrast
would delimit cross-family replication.  The complete cluster table is
retained regardless of direction.

All six Granite cluster effects are positive.  The separately retained Qwen
contrast is also positive in all six clusters, with a mean field-minus-nominal
peak of 112.691 units.  These absolute values are not compared across models:
each leave-one-cluster-out geometry is standardized and fitted within model.

\resultfigure{figure06_cross_model_quench.pdf}{0.94\textwidth}{%
Model-specific field-quench and nominal trajectories.  Curves are means over
six graph/environment clusters and bands show between-cluster dispersion.
Nominal fits, scaling, covariance, and thresholds exclude the displayed
cluster.}{fig:crossquench}

Restoration H4 is \csname V15HFour\endcsname{} distance units (95\% interval
\csname V15HFourCI\endcsname), with frozen disposition
\csname V15HFourDisposition\endcsname{}.  Unlike the historical V14 H3
statistic, this fixed-window quantity permits either sign.  Recovery time and
threshold crossing remain trajectory descriptors, not an equilibrium
relaxation theorem.  The quench and counter-quench paths need not retrace one
another because messages, actions, and history alter the state reached at
restoration.

The fixed-window decline occurs in every model--cluster pair, but it does not
imply complete return within the observation horizon.  Every Qwen trajectory
re-enters its training-nominal threshold six sweeps after restoration.  Only
one of six Granite trajectories re-enters by sweep 45; the other five retain
final-five means above their own thresholds.  Thus H4 supports motion toward
the restored nominal regime, while Granite exhibits slower or incomplete
finite-horizon recovery.

\section{Results: memory and temporal asymmetry}

Relative to Markovized state, genuine memory changes adjusted block divergence
by \csname V15HTwo\endcsname{} nats per attempted update (95\% interval
\csname V15HTwoCI\endcsname); the frozen disposition is
\csname V15HTwoDisposition\endcsname{}.  Relative to the scrambled-history
placebo, the change is \csname V15HThree\endcsname{}
(\csname V15HThreeCI\endcsname), disposition
\csname V15HThreeDisposition\endcsname{}.  H2 asks whether history matters
relative to an
explicit current-state prompt.  H3 is more stringent: it asks whether the
temporally related content matters beyond adding a similarly shaped prompt
section.

Several arm-level adjusted values are negative because their observed block
divergence lies below the corresponding finite-sample shuffle-floor mean.  H2
and H3 are paired differences between these untruncated adjusted statistics;
their positive contrasts do not establish positive absolute entropy production
in either arm.

The prospectively pooled contrast retains substantial model heterogeneity.
For persistent minus Markovized state, model-specific means are 0.03041 for
Qwen (five of six cluster effects positive) and 0.07835 for Granite (six of
six positive).  For persistent minus scrambled history they are 0.00689 for
Qwen (four of six positive) and 0.09001 for Granite (six of six positive).
The model-specific Qwen scrambled-control sign-flip value is 0.21875 and is
not separately confirmatory; the frozen pooled H3 remains supported because
model-stratified cluster pairs were the prespecified units.  This decomposition
supports cross-family direction for the Markovized contrast but not a claim
that content-specific memory effects are uniform across model families.

\resultfigure{figure07_v15_memory_controls.pdf}{0.94\textwidth}{%
Per-cluster genuine-memory contrasts for both pinned model families.  Short
horizontal segments indicate model means.  The scrambled-history arm controls
prompt section and approximate length while destroying temporal relation to
the current trajectory.}{fig:v15memory}

No exact thermodynamic interpretation follows from a positive result.  The
observed block words omit natural-language content and some private state.
Conversely, a null H3 would be informative: prompt presence, bounded memory
depth, model-specific history use, or estimator power could explain the V12
and V13 contrast without a reproducible content-specific effect.  Block-length,
pseudocount, conditional-memory-depth, occupancy, and token-balance tables
support this distinction.

Estimator sensitivity further limits the scale claim.  The pooled memory
contrasts remain positive over most prespecified block-length and pseudocount
combinations, but their magnitude varies substantially; at block length four
and pseudocount one, persistent-minus-scrambled is only 0.00024 nats per
update and the Granite component is slightly negative.  Conditional-memory
information is likewise not monotonically larger in the persistent arm.
Accordingly, the robust statement concerns the frozen coarse graining, its
bias floor, and paired controls rather than an estimator-independent entropy-
production rate.

\section{Spatial correlation and finite-window persistence}

Equations (\ref{eq:connected-correlation})--(\ref{eq:binder}) are secondary
descriptive analyses of the same frozen trajectories; they neither modify the
confirmatory estimands nor introduce additional simulation arms.  They were
specified during reconstruction after the original V15 aggregate results were
known and are therefore not a new prospective hypothesis family.  Connected
correlation matrices are first calculated within every complete
model--graph--environment trajectory and phase.  Their graph-distance profiles
are then summarized over the six independent clusters.  Figure
\ref{fig:correlations} shows cluster-averaged matrices for the persistent-
history quench window and phase-resolved distance profiles.  Community
boundaries are fixed by the graph construction rather than selected from the
observed correlations.

At graph distance one, the persistent-history quench-minus-baseline connected
correlation changes are \csname V15QwenCorrQuench\endcsname{}
(\csname V15QwenCorrQuenchCI\endcsname) for Qwen and
\csname V15GraniteCorrQuench\endcsname{}
(\csname V15GraniteCorrQuenchCI\endcsname) for Granite.  These cluster-level
descriptive intervals quantify finite-window reorganization; their signs are
not assigned a confirmatory disposition.

\resultfigure{figure13_graph_distance_correlations.pdf}{0.94\textwidth}{%
Connected belief correlations in the persistent-history arm.  The upper
panels average complete $16\times16$ connected-correlation matrices over all
six clusters in each model; black rules mark the two predefined communities.
Agent indices align across graph realizations only through those community
memberships, so fine node-level texture in an averaged matrix has no
cross-realization identity interpretation.
The lower panels show equation (\ref{eq:connected-correlation}) by graph
distance, with 95\% complete-cluster bootstrap intervals.  Individual pairs
and updates do not enter the inferential sample size.}{fig:correlations}

Supplementary figure S1 combines the full autocorrelation curve with the
fixed-lag integral in equation (\ref{eq:autocorrelation}), the Binder statistic
in equation (\ref{eq:binder}), and empirical magnetization occupancy.  This
combination prevents a Binder value from being interpreted without its
underlying finite-window distribution.  It also separates persistence of the
projected collective coordinate from the block-reversal statistic: the former
is time-symmetric two-point dependence, whereas the latter compares forward
and reversed path words.
The accompanying source table also retains full-, early-half-, and late-half
Binder estimates under cluster-first and pooled-moment rules.  Their
differences diagnose finite-window and aggregation sensitivity rather than
selecting whichever construction yields the larger separation.

At the fixed two-sweep truncation, persistent-minus-Markovized recovery
$\tau_{\mathrm{int}}$ changes are \csname V15QwenTauMemory\endcsname{}
(\csname V15QwenTauMemoryCI\endcsname) updates for Qwen and
\csname V15GraniteTauMemory\endcsname{}
(\csname V15GraniteTauMemoryCI\endcsname) updates for Granite.  The
field-Markovized quench-minus-baseline Binder changes are
\csname V15QwenBinderQuench\endcsname{}
(\csname V15QwenBinderQuenchCI\endcsname) and
\csname V15GraniteBinderQuench\endcsname{}
(\csname V15GraniteBinderQuenchCI\endcsname), respectively.  These estimates
are reported whether positive, negative, or interval-compatible with zero.

\section{Out-of-sample kinetic closure}

We fit logistic belief and action responses only to the immutable V13
microscopic-response table.  The belief model includes private field, delivered
neighbor field times controlled coupling, previous belief, and previous action;
the action model includes updated belief and previous action.  No coefficient
is refitted to a V14 or V15 quench trajectory.  The surrogate is run on the
same graph generator, initial-state generator, asynchronous update tape, field
schedule, and intervention times.  It is a deliberately low-dimensional
kinetic-Ising comparison rather than a fitted description of every asymmetric
or history-dependent response \cite{aguilera2021}.

\resultfigure{figure10_direct_surrogate_quench.pdf}{0.94\textwidth}{%
Out-of-sample V14 field-quench comparison between direct Qwen trajectories and
the V13-fitted kinetic surrogate.  Matching direction, peak timing, normalized
amplitude, recovery, and route area are assessed separately.  Surrogate
failure is evidence about the limits of the closure, not a failed LLM
experiment.}{fig:surrogate}

In the five-coordinate comparison geometry, the direct field-minus-nominal
peak-response difference is only 0.142 units on average and is nonzero in one
of six clusters, whereas the surrogate difference is 1.467 units and is
positive in four clusters.  The corresponding field-minus-nominal
energy--entropy route-area differences are 0.043 and 0.695.  The direct shared
response peaks at the quench boundary in one cluster and is otherwise dominated
by initial relaxation; the surrogate places its peak at the quench or
counter-quench boundary in four clusters.  Both descriptions return close to
their own final shared-response baseline, but the surrogate neither reproduces
the direct amplitude nor its timing.  This also explains why the richer V14
macrostate distance can show a large quench pulse that a short kinetic closure
over a small shared coordinate set misses.

This comparison identifies three possible failure modes.  First, a homogeneous
logistic neighbor coefficient cannot capture state-dependent semantic use.
Second, action and belief persistence can introduce lags not represented by a
single synchronous mean field.  Third, hidden memory and workload can produce
different future laws at the same projected spin configuration.  Adding
unconstrained parameters after observing quench trajectories would obscure
these limits, so the closure remains low-dimensional and out of sample.

CPU trajectories at $N\in\{8,16,32,64\}$ supply a dense effective-model size
context.  They are not direct LLM finite-size scaling and cannot establish a
critical exponent or thermodynamic-limit phase transition.  The direct
$N=16$ paths remain empirical anchors.

\section{What the reduced representation contributes}

Magnetization records directional order but can remain near zero when half the
population reorganizes in opposing communities.  Configuration entropy then
records collective state diversity; marginal entropy records local occupancy;
total correlation distinguishes independent uncertainty from coordinated
dependence; effective energy records compatibility with the symmetric
reference; susceptibility and energy variance record fluctuations; correlation
time records persistence; and path reversal records temporal ordering.  These
quantities are complementary only when their estimators and scales are kept
distinct.

The V14 representation experiment uses transparent multinomial logistic
regression with leave-one-cluster-out preprocessing.  The delayed permutation
audit tests both absolute balanced accuracy and the increment over order-only
features.  This is not a claim that entropy outperforms all ordinary summaries.
It tests whether a predefined collective representation retains condition
information discarded by a reduced representation.  The field-quench
trajectory supplies a mechanistic example: energy, dependence, entropy rate,
and fluctuation can change while mean magnetization remains balanced.

\resultfigure{figure09_confirmatory_effects.pdf}{0.94\textwidth}{%
V15 frozen effects and 95\% cluster-bootstrap intervals.  Distance and
information units occupy separate panels; no claims forest places incompatible
quantities on one numerical axis.  Dispositions use exact cluster sign flips
and the frozen multiplicity allocation.}{fig:effects}

\section{Interpretation and limitations}

The formulation has a literal stochastic-process component and an effective
statistical-mechanics component.  Agent instances, local prompts, packet
delivery, update schedules, and categorical state changes are implemented
objects.  Magnetization, Shannon entropy, mutual information, total
correlation, and path KL are mathematical observables of recorded data.
Reference energy is effective because the LLM transition rule is not derived
from its Hamiltonian.  Decoding temperature is a model-generation setting.
No free-energy-like coordinate is emphasized because an independently
identified effective temperature is unavailable.

Several limitations bound the claim.  First, two 7B instruction-model families
do not establish universality.  Second, $N=16$ direct systems are finite; no
thermodynamic-limit transition is claimed.  Third, the binary projection
compresses semantic messages and private state, so temporal asymmetry is a
coarse-grained diagnostic.  Fourth, only a reciprocal modular topology and one
coupling/noise anchor are used in V15.  Fifth, a scrambled-history placebo
matches format and approximate token distribution but cannot make semantic
content exchangeable.  Sixth, short-window entropy and dependence estimates
remain coarse-graining dependent even after a null-floor audit.

The study also preserves negative boundaries.  V12 found no reliable increase
in irreversibility from degree- and traffic-matched nonreciprocity.  V14
partition and message-corruption trajectories remained close to nominal under
their tested construction.  A simple kinetic closure missed parts of the
direct coupling response.  These limits prevent the broad assertion that any
directed graph, entropy statistic, or Ising analogy is automatically useful.

\section{Conclusion}

State-separated LLM-agent instances form a measurable interacting stochastic
system whose collective trajectories admit a disciplined reduced description.
The distinction between complete augmented state, observable projection, and
rolling macrostate explains why private memory can generate projected history
dependence.  Controlled field quench and restoration create a common test of
order, uncertainty, dependence, effective compatibility, fluctuation, and
temporal response.  V15 adds an independent model family and a format-matched
history control to the V12--V14 evidence.  The pooled prospective memory
effects are positive, while the weaker Qwen scrambled-control decomposition,
Granite's incomplete threshold re-entry, and estimator sensitivity remain
explicit boundary results.

The statistical-mechanical contribution is therefore not that entropy makes
agents perform better.  It is a formulation and measurement framework that
connects local language-conditioned transitions to collective finite-size
dynamics, exposes what population averages discard, and makes failures of
simple closures explicit.  Further work requires more model families, direct
sizes and topologies, longer stationary trajectories, and observable-state
augmentation before thermodynamic-limit or exact stochastic-thermodynamic
claims would be warranted.

\section*{Reproducibility and data availability}

The repository records immutable model revisions, package versions, prompt and
schema hashes, graph and seed manifests, control-tape hashes, frozen protocol,
analysis configuration, aggregate tables, figure source data, vector PDFs, and
content-addressed manifests.  Raw prompts, completions, and full trajectories
remain outside Git because they are large and potentially sensitive; their
tree hashes and replay results are repository-facing.  Every trajectory can be
regenerated through content-addressed model decisions.  The original sealed
V15 execution used \csname V15OriginalGPUHours\endcsname{} measured generation
GPU-hours; the fresh from-scratch reconstruction accounted for at least
\csname V15ReconstructionGPUHours\endcsname{} measured generation GPU-hours
under the same scientific protocol.  This runtime difference is an
environment-dependent reproducibility cost rather than a scientific outcome.
The original external raw tree did not survive a later Pod replacement.  The
fresh call records replay internally at decision resolution and their frozen
aggregate tables, calls, and token counts are compared with the committed
reference package; historical call-file digest identity cannot be checked
without the deleted original raw tree.
The content-addressed accounting audit separately identified
\csname V15OrphanFormalDecisions\endcsname{} decision records
(\csname V15OrphanFormalGPUHours\endcsname{} generation GPU-hours) from any
interrupted atomic panels; these records contribute to compute accounting but
not to a scientific trajectory.
One additional model generation failed after inference but before its atomic
raw record could be written.  It is counted as
\csname V15UnrecordedInfrastructureCalls\endcsname{} infrastructure call, but
its tokens and latency are unavailable; fresh reconstruction token, GPU-hour,
and cost totals are therefore measured lower bounds.
No chain of
thought was requested or retained.

Every experimental arm instantiates a fresh agent-network state from its
frozen panel seeds.  Model weights and the tokenizer are shared read-only for
throughput, but conversational history, generation key/value caches, mutable
agent objects, and scientific random-number state are not carried between
arms.  Each generation resets the frozen per-decision seed; record counters
and provider accounting do not enter prompts or transition probabilities.

The frozen execution-source digest inadvertently enumerated one ignored
224-byte Python bytecode cache outside a conventional \texttt{\_\_pycache\_\_}
directory.  The cache was not scientific source and is not distributed as a
repository artifact.  A versioned provenance audit reconstructs the legacy
digest from its recorded path and hash and additionally reports a reproducible
cache-free semantic-source digest.  This correction changes neither the
protocol nor any formal trajectory or result.

The manuscript uses standard 12-point LaTeX, which current IOP guidance accepts
alongside the optional \texttt{iopjournal} class.  Figures are embedded vector
PDFs with source CSVs and are rendered at 300 DPI for inspection.

\section*{AI-assisted preparation declaration}

AI-assisted software and text preparation was used under human-directed
protocol constraints.  All scientific estimands, code paths, numerical tables,
citations, figures, and claims remain subject to author verification.  No
language-model output is treated as human evidence.

\appendix
\section{Frozen inferential details}

For $n\leq20$ paired clusters, the one-sided sign-flip probability enumerates
all $2^n$ sign assignments of the observed cluster effects.  The cluster
bootstrap samples complete graph/environment clusters with replacement and
never resamples time points as independent observations.  H1 receives alpha
0.02; H2--H4 receive a Holm-adjusted family alpha 0.03.  The allocation was
fixed before formal outcomes.  Intervals are uncertainty summaries and do not
replace the exact paired tests.

The V14 representation permutation retains the four condition panels in every
cluster and permutes labels only within cluster.  For each of 10,000
replicates, every leave-one-cluster-out fold refits imputation, scaling, and the
multinomial logistic model.  Test-cluster features do not enter preprocessing.

\section{Estimator sensitivity}

Total-correlation nulls independently circular-shift each coordinate by a
nonzero offset where possible.  This retains marginal occupancy and each
coordinate's internal temporal ordering while breaking synchronous dependence.
The adjusted estimate is not truncated.  Alternative rolling windows rebuild
the words, covariance, training threshold, and held-out distance rather than
subsampling a five-sweep output.

For block reversal, time-shuffled floors quantify finite-length occupancy
bias.  They do not prove stationarity.  Results at block lengths two, three,
and four expose finite-history sensitivity; conditional mutual information at
history depths one through three provides a separate diagnostic of incomplete
first-order closure.

\section{Authority and privacy checks}

An agent's response is causally consequential: it changes the scheduled
instance's recorded belief, action, commitment, memory code, signal, and typed
tool consequence.  The scheduler validates but does not replace those fields.
Every transition records the scheduled identity, recipient, packet bytes,
prompt hash, output-record hash, and fingerprints of unrelated peer-private
state.  Tests perturb one instance's memory or decision provider while
confirming that peer-private stores remain unchanged.

\bibliographystyle{unsrtnat}
\bibliography{references}

\end{document}
